\documentclass[10pt,aspectratio=169]{beamer}

\usetheme{Yantian}

\usepackage{lipsum}
%\providecommand{\lishu}{\CJKfontspec{Libian SC}[BoldFont=Baoli SC,Scale=1.2]}
%
%\setbeamertemplate{sections/subsections in toc}[ball]
\usepackage{amsmath}
\usepackage{xeCJK}
\usepackage{caption}
%\usepackage{mathspec}
%\setmainfont{Times New Roman}
%\setmathsfont{Times New Roman}
\setCJKmainfont{KaiTi}[AutoFakeBold] %FZKTK--GBK1-0
\usefonttheme{professionalfonts}
\def\mathfamilydefault{\rmdefault}
\DeclareMathOperator*{\argmin}{arg\,min}
\DeclareMathOperator*{\argmax}{arg\,max}
\usepackage{bm}
\usepackage{multicol}
\usepackage{multirow}
\usepackage{booktabs}
\usepackage{colortbl}
\usepackage{bbm}
\usepackage{extarrows}
\usepackage{pgfornament-han}
%\usepackage{enumitem}
\usepackage{fontawesome5}
\usepackage{subfigure}
\usetikzlibrary{chains}
\usetikzlibrary{calc}
\usetikzlibrary{decorations}
\usetikzlibrary{decorations.text,decorations.markings}
\usetikzlibrary{fadings}
\renewcommand{\arraystretch}{1.5}
\definecolor{secolor}    {RGB}{176, 164, 227}
%\setbeamertemplate{sections/subsections in toc}[ball]
\makeatletter
\patchcmd{\beamer@sectionintoc}
{\vfill}
{\vskip 10pt}
{}
{}
\makeatother


\title{Template}
%\subtitle{子标题}
\author[罗雁天]{罗雁天}
\mail{lyt18@mails.tsinghua.edu.cn}
\institute[清华大学]{无线多媒体通信实验室\\
清华大学电子工程系}
\date{\today}
\renewcommand{\today}{\number\year 年 \number\month 月 \number\day 日}

\AtBeginSection[]{
	{\usebackgroundtemplate{\includegraphics[width=\paperwidth]{bg/toc.png}}
	\begin{frame}[plain]
		\begin{columns}
			\column{.2\linewidth}
			\begin{tikzpicture}[remember picture,overlay]
				\node at (current page.north west) [scale=2.5,xshift=1.7cm,yshift=-.35cm,font=\chieu \lishu \bfseries,text=DarkPurple]{目录};
			\end{tikzpicture}
			\column{.8\linewidth}
			\chieu \lishu \bfseries \Large
			\tableofcontents[currentsection]
		\end{columns}
	\end{frame}}
}

%\AtBeginSubsection[]{
%	\begin{frame}[plain]{OUTLINE}
%		\begin{ascolorbox12}{OUTLINE}
%			\chieu \lishu \bfseries
%			\tableofcontents[currentsection,currentsubsection]
%		\end{ascolorbox12}
%	\end{frame}
%}


\begin{document}

\maketitle

{
	\usebackgroundtemplate{\includegraphics[width=\paperwidth]{bg/toc.png}}
	\begin{frame}[plain]
		\begin{columns}
			\column{.2\linewidth}
			\begin{tikzpicture}[remember picture,overlay]
				\node at (current page.north west) [scale=2.5,xshift=1.7cm,yshift=-.35cm,font=\chieu \lishu \bfseries,text=DarkPurple]{目录};
			\end{tikzpicture}
			\column{.8\linewidth}
			\chieu \lishu \bfseries \Large
			\tableofcontents
		\end{columns}
	\end{frame}
}

\section{Introduction}
\begin{frame}{Introduction}
		\begin{itemize}
			\item 修改字体命令在beamerthemeYantian.sty里，Mac上的字体和Windows上的字体不一样，需要手动修改
			\item 志明是雁天儿子
		\end{itemize}
\end{frame}

\section{乱七八糟凑字数}
\begin{frame}{乱七八糟凑字数}
	\begin{itemize}
		\item 关山难越，谁悲失路之人；萍水相逢，尽是他乡之客
		\item 我达达的马蹄是美丽的错误，我不是归人，是个过客
		\item 不为盛名而来，不为低谷而去
	\end{itemize}
\end{frame}

{
	\usebackgroundtemplate{\begin{tikzpicture}[overlay, remember picture]
			\node[opacity=1,xshift=8.0cm,yshift=-4.49cm] {\includegraphics[width=\paperwidth,height=\paperheight]{bg/bg1}};
		\end{tikzpicture}}
	\begin{frame}[plain]
		\begin{center}
			\begin{tikzpicture}[x=1pt,y=1pt]
			\node[font=\Huge \huati \lishu \bfseries,align=center,xshift=120,yshift=-25,text=white] at (0,0) {Thank You!};
			\node[font=\Huge \huati \lishu \bfseries,align=center,xshift=120,yshift=-45,text=white,yscale=-1,scope fading=south, opacity=0.5] at (0,0) {Thank You!};
			\end{tikzpicture} 
		\end{center}
	\end{frame}
}

\end{document}


